\documentclass[12pt, titlepage]{article}
\usepackage{amsmath}
\usepackage{amsfonts}
\usepackage{amssymb}
\usepackage{fancyhdr}
\usepackage{cancel}
\usepackage{xcolor}
\usepackage[bidi=basic]{babel}
\babelprovide[main,import]{hebrew}
\babelfont[hebrew]{rm}{Hadasim CLM}
\babelfont[hebrew]{sf}{Miriam Mono CLM}
\usepackage{titlesec}
\newcommand\Ccancel[2][black]{\renewcommand\CancelColor{\color{#1}}\cancel{#2}}
\newcommand{\pard}[2][f]{\frac{\partial #1}{\partial #2}}
\title{חשבון אינפי 2 - 01040281 \\ גליון 6}
\author{יונתן אבידור - 214269565}
\titleformat{\section}{\Large\sffamily\bfseries}{\thesection}{1em}{}
\begin{document}
\maketitle
\pagestyle{fancy}
\fancyhead{}
\fancyfoot[L]{\text{214269565}}
\part*{שאלה 1}
בדקו האם הגבולות הבאים קיימים
\section*{סעיף א'}
\[
    \lim_{\left(x,y\right)\to\left(0,0\right)} \frac{1-\cos 2xy}{x^2y\sin \left(\pi y\right)}
\]
\section*{סעיף ב'}
\[
    \lim_{\left(x,y\right)\to\left(0,0\right)} \frac{x^4-x^2-y^4+y^2}{x+y}
\]
\part*{פתרון 1}
\section*{סעיף א'}
\section*{סעיף ב'}
\newpage
\part*{שאלה 2}
תהא
\[
    f\left(x,y\right) = \begin{cases}
        \frac{x^3+y^3}{2x^2+y^2} & \left(x,y\right) \neq(0,0)\\
        0 & \left(x,y\right) = \left(0,0\right)
    \end{cases}
\]
האם $f$ רציפה ב-$\left(0,0\right)$? האם $\frac{\partial f}{\partial x} \left(0,0\right)$ ו-$\frac{\partial f}{\partial x} \left(0,0\right)$ קיימים? האם $f$ דיפרנציאבילית ב-$\left(0,0\right)$
\part*{פתרון 2}
\newpage
\part*{שאלה 3}
תהא
\[
    f\left(x,y\right) = \begin{cases}
    \frac{x^2y}{\left(x^2+y^2\right)^k} & \left(x,y\right) \neq \left(0,0\right)\\
        0 & \left(x,y\right) = \left(0,0\right)
    \end{cases}
\]
\section*{סעיף א'}
עבור אילו $k>0$ $f$ רציפה ב-$\left(0,0\right)$?
\section*{סעיף ב'}
בדקו האם הנגזרות החלקיות קיימות ב-$\left(0,0\right)$, אם כן. חשבו אותם
\section*{סעיף ג'}
עבור אילו $k>0$ $f$ דיפרנציאבילית ב-$\left(0,0\right)$?
\part*{פתרון 3}
\section*{סעיף א'}
נציב
\begin{gather*}
    x = r \cos \theta \\
    y = r \sin \theta
\end{gather*}
כעת נבדוק את הגבול
\[
    \lim_{\left(x,y\right) \to \left(0,0\right)} f\left(r\cos \theta,r\sin\theta\right)
\]
אם הגבול הוא $0$, הפונקציה רציפה.
\begin{gather*}
    f\left(r\cos \theta, r\sin\theta\right) = \frac{r^2 \cos^2\theta r\sin\theta}{\left(r^2\cos^2\theta+r^2\sin^2\right)^k} = \frac{r^3 \cos^2\theta \sin\theta}{\left(r^2\underbrace{\left(\cos^2\theta+\sin^2\theta\right)}_{=1}\right)^k}\\
    =\frac{r^3 \cos^2\theta \sin\theta}{\left(r^2\right)^k} = \frac{r^3 \cos^2\theta \sin\theta}{r^{2k}} = \frac{r^3}{r^{2k}} \cos^2\theta\sin\theta=r^{3-2k} \cos^2\theta\sin\theta
\end{gather*}
$\cos^2\theta\sin\theta$ חסומה, לכן אם $r^{3-2k}\xrightarrow{r\to0}0$ יתקיים $f\left(x,y\right)\xrightarrow{\left(x,y\right)\to\left(0,0\right)}0$
כאשר $r\to 0$, נרצה שיתקיים $3-2k>0$, כי אם $3-2k=0$  יהיה $r^0$ שזה $1$, ואם יתקיים $3-2k<0$ אז $r^{3-2k}\xrightarrow{r\to0}\infty$
\[
3-2k > 0 \implies 3 > 2k \implies \boxed{\frac{3}{2} > k}
\]
\section*{סעיף ב'}
\[
    \frac{\partial f}{\partial x} \left(0,0\right) = \lim_{\Delta x} \frac{\frac{\left(0+\Delta x\right)^2 \cdot 0}{\left(\left(0+\Delta x\right)^2 + 0^2\right)^k} - 0}{\Delta x} = 0
\]

\[
    \frac{\partial f}{\partial y} \left(0,0\right) = \lim_{\Delta y} \frac{\frac{0^2 \cdot \left(0+\Delta y\right)}{\left(0^2 + \left(0+\Delta y\right)^2\right)^k} - 0}{\Delta y} = 0
\]
\section*{סעיף ג'}
הראינו כי $f$ רציפה ב-$\left(0,0\right)$ ושתי הנגזרות קיימות בו, לכן הפונקציה דיפרנציאבילית אם
\[
    \lim_{\left(x,y\right) \to \left(0,0\right)} \frac{f\left(x,y\right) -f\left(0,0\right) - \frac{\partial f}{\partial x} \left(0,0\right) \cdot \left(x-0\right) - \frac{\partial f}{\partial y} \left(0,0\right) \cdot \left(y-0\right)}{\sqrt{\left(x-0\right)^2 + \left(y-0\right)^2}} = 0
\]
בסעיף הוקדם הראינו כי $\frac{\partial f}{\partial y} = 0, \frac{\partial f}{\partial x} = 0$
לכן אנחנו בעצם מחפשים את הגבול הבא
\begin{gather*}
    \lim_{\left(x,y\right) \to \left(0,0\right)} \frac{f\left(x,y\right)}{\sqrt{x^2 +y^2}} = \lim_{\left(x,y\right) \to \left(0,0\right)} \frac{x^2y}{\left(x^2+y^2\right)^k\cdot \left(x^2+y^2\right)^{\frac{1}{2}}} \\
    = \lim_{\left(x,y\right) \to \left(0,0\right)} \frac{x^2y}{\left(x^2+y^2\right)^{k+ \frac{1}{2}}}
\end{gather*}
בסעיף א' ראינו כי הגבול הוא $0$ אמ"מ $3-2\left(k+\frac{1}{2}\right)>0$
\[
3-2\left(k+\frac{1}{2}\right) > 0 \implies 3 -2k-1 > 0 \implies 2 > 2k \implies  \boxed{1 > k}
\]
\newpage
\part*{שאלה 4}
חשבו את $f_{xy}$ ואת $f_{yx}$ ב-$\left(0,0\right)$ כאשר
\[
    f(x,y) = \begin{cases}
        x^2 \sin \frac{1}{x} + y^2 \sin \frac{1}{y} & x,y \neq 0\\
        x^2 \sin \frac{1}{x} & x \neq 0 ,y = 0\\
        y^2 \sin \frac{1}{y} & x = 0, y \neq 0\\
        0 & x=y=0
    \end{cases}
\]
\part*{פתרון 4}
\newpage
\part*{שאלה 5}
תהא $f\left(x,y\right)$ המקיימת $\left|f\left(x,y\right)\right| \leq M\left(x^2+y^2\right)$ עבור $M>0$ כלשהו ולכל $\left(x,y\right)$ בסביבה של $\left(0,0\right)$. הוכיחו כי $f$ דיפרנציאבילית ב-$\left(0,0\right)$
\part*{פתרון 5}
\newpage
\part*{שאלה 6}
תהא $f$ דיפרנציאבילית ב-$\left(x_0,y_0\right)$. הוכיחו כי קיים כיוון $v$ שעבורו הנגזרת המכוונת של $f$ היא $0$
\part*{תשובה 6}
נחפש $v = \begin{pmatrix}
    v_1,v_2
\end{pmatrix} \in \mathbb{R}^2$ המקיים
\[
  \pard{v} \left(x_0,y_0\right) = \left<\nabla f\left(x_0,y_0\right), v \right> = 0
\]
לכל שני וקטורים $u,w$ מתקיים
\[
    \left<u,w\right> = 0 \iff u \perp w
\]
לכן נדרוש ש $\left(v_1,v_2\right) \perp \left(x_0,y_0\right)$\\
נייצג את $\left(x_0,y_0\right)$ בצורה פולארית
\[
\begin{gathered}
    x_0 = r \cos \theta\\
    y_0 = r \sin \theta
\end{gathered} \implies \left(x_0,y_0\right) = \left(r\cos\theta,r\sin\theta\right)
\]
על מנת למצוא וקטור ניצב, נזיז את $\left(x_0,y_0\right)$ בתשעים מעלות.
\[
    \left(v_1,v_2\right) = \left(r\cos\left(\theta + \frac{\pi}{2}\right),r\sin\left(\theta+\frac{\pi}{2}\right)\right) = \left(-r \sin\theta, r\cos\theta\right) = \left(-y_0,x_0\right)
\]
לכן
\[
    \boxed{v =\left( 
            -y_0,x_0
    \right)}
\]
נחשב ונבדוק
\[
    \left<\nabla f\left(x_0,y_0\right), \left(-y_0,x_0\right)\right>
\]
\newpage
\part*{שאלה 7}
מצאו את משוואת המישור המשיק ואת הנורמל של המישור המשיק של הפונקציה $z=e^{x+\sin y}$ בנקודה $\left(0,0,1\right)$
\part*{תשובה 7}
\newpage
\part*{שאלה 8}
יהי $m \in \mathbb{R}$. הפונקציה $f\left(x,y\right)$ נקראית הומוגנית מסדר $m$ אם היא מקיימת 
\[
    f\left(\lambda x, \lambda y\right) = \lambda^m f\left(x,y\right)
\]
לכל $x,y,\lambda$
\section*{סעיף א'}
הראו שכל פונקציה הומוגנית מסדר $m$ יכולה להיות מיוצגת בתור 
\[
    f\left(x,y\right) = x^m F\left(\frac{y}{x}\right)
\]
לכל $x\neq 0$, כאשר $F$ היא פונקציה כלשהי במשתנה אחד
\section*{סעיף ב'}
הראו שאם $f$ הומוגנית מסדר $m$ אז היא מקיימת את המשוואה הדיפרנציאלית
\[
    x f_x + y f_y = m f
\]
\part*{תשובה 8}
\section*{סעיף א'}
\section*{סעיף ב'}
\end{document}
